\documentclass{article}

\usepackage[italian]{babel}
\usepackage[margin=2cm, footskip=5mm]{geometry}
% questi package non sono necessari in lualatex; ref https://tex.stackexchange.com/a/413046
% \usepackage[utf8]{inputenc}
% \usepackage[T1]{fontenc}
\usepackage{enumitem}
\usepackage{hyperref}
\usepackage{titlesec}
\usepackage{soulutf8}
\usepackage{contour}
\usepackage{float}
\usepackage{graphicx}
\usepackage{fancyhdr}
\usepackage{longtable}
\usepackage[table]{xcolor}
\usepackage{titling}
\usepackage{lastpage}
\usepackage{ifthen}
\usepackage{calc}
\usepackage{minted}
\usepackage{pgfgantt}
\usepackage{subfiles}

\newlength{\imgwidth}

\newcommand\scalegraphics[1]{%
    \settowidth{\imgwidth}{\includegraphics{#1}}%
    \setlength{\imgwidth}{\minof{\imgwidth}{\textwidth}}%
    \includegraphics[width=\imgwidth]{#1}%
}

% XXX definizione dei percorsi in cui cercare immagini
\graphicspath{ {./}
    {./img/}
}

% esempio di utilizzo: \appendToGraphicspath{./img/} (un comando diverso per ogni path da includere)
% N.B.: ci DEVE essere un forward slash alla fine del path, a indicare che è una cartella.
\makeatletter
\newcommand\appendToGraphicspath[1]{%
  \g@addto@macro\Ginput@path{{#1}}%
}
\makeatother

% setup della sottolineatura
\setuldepth{Flat}
\contourlength{0.8pt}

\newcommand{\uline}[1]{%
  \ul{{\phantom{#1}}}%
  \llap{\contour{white}{#1}}%
}

% setup dei link
\hypersetup{
  colorlinks=true, % set true if you want colored links
  linktoc=all,     % set to all if you want both sections and subsections linked
  linkcolor=black, % choose some color if you want links to stand out
}

% setup di header e footer
\pagestyle{fancy}

\fancyhf{}
\fancyhead[L]{\includegraphics[width=1cm]{logo.png}}
\fancyhead[R]{\thetitle}
\fancyfoot[R]{\thepage\ di~\pageref{LastPage}}

\fancypagestyle{nopage}{%
  \fancyfoot{}%
}

\setlength{\headheight}{1.2cm}

% setup forma \paragraph e \subparagraph
\titleformat{\paragraph}[hang]{\normalfont\normalsize\bfseries}{\theparagraph}{1em}{}
\titleformat{\subparagraph}[hang]{\normalfont\normalsize\bfseries}{\thesubparagraph}{1em}{}

% setup profondità indice di default
\setcounter{secnumdepth}{5}
\setcounter{tocdepth}{5}

% shortcut per i placeholder
\newcommand{\plchold}[1]{\textit{\{#1\}}} % chktex 20

% hook per lo script che genera il glossario
\newcommand{\glossario}[1]{\underline{#1}\textsubscript{g}}

% definizione dei comandi \uso e \stato
\makeatletter
\newcommand{\setUso}[1]{%
  \newcommand{\@uso}{#1}%
}
\newcommand{\uso}{\@uso}

\newcommand{\setStato}[1]{%
  \newcommand{\@stato}{#1}%
}
\newcommand{\stato}{\@stato}

\newcommand{\setVersione}[1]{%
  \newcommand{\@versione}{#1}%
}
\newcommand{\versione}{\@versione}

\newcommand{\setResponsabile}[1]{%
  \newcommand{\@responsabile}{#1}%
}
\newcommand{\responsabile}{\@responsabile}

\newcommand{\setRedattori}[1]{%
  \newcommand{\@redattori}{#1}%
}
\newcommand{\redattori}{\@redattori}

\newcommand{\setVerificatori}[1]{%
  \newcommand{\@verificatori}{#1}%
}
\newcommand{\verificatori}{\@verificatori}

\newcommand{\setDescrizione}[1]{%
  \newcommand{\@descrizione}{#1}%
}
\newcommand{\descrizione}{\@descrizione}

\newcommand{\setModifiche}[1]{%
  \newcommand{\@modifiche}{#1}%
}

\newcommand{\modifiche}{\@modifiche}

\makeatother

% setup delle description
\setlist[description,1]{font=$\bullet$\hspace{1.5mm}, labelwidth=* leftmargin=*,labelindent=12.5mm}
\setlist[description,2]{font=$\bullet$\hspace{1.5mm}, leftmargin=*,labelindent=12.5mm}

\appendToGraphicspath{../../commons/img/}

\title{Verbale interno --- 20/04/2020}

\setResponsabile{Luca Ercole}
\setRedattori{Alberto Cocco}
\setVerificatori{
  Alberto Gobbo
}
\setUso{Interno}
\setDescrizione{Verbale dell'incontro di GruppOne del 20/04/2020}
\setModifiche{%
\cellcolor{white!80!lightgray!100} & Luca Ercole & 2020--04--21 & approva documento \\%
\cellcolor{white!80!lightgray!100} & Alberto Gobbo & 2020--03--20 & verifica verbale \\%
\multirow{-3}{*}{-} & Alberto Cocco & 2020--04--20 & stendi verbale%
}

\disabilitaVersione{}
\disabilitaElencoFigure{}
\disabilitaElencoTabelle{}

\begin{document}

\thispagestyle{empty}
\pagenumbering{gobble}

\begin{center}

  \includegraphics[width=8.5cm]{\commons/img/logo.png}\\
  {\Large GruppOne - progetto "Stalker"}\\
  \vspace{1.5cm}

  {\Huge \thetitle}
  \vspace{1.5cm}

  \begin{table}[H]
    \centering

    \begin{tabular}{r|l}
      \textbf{Versione}     & \versione              \\
      \textbf{Approvazione} & \responsabile          \\
      \textbf{Redazione}    & \redattori             \\
      \textbf{Verifica}     & \verificatori          \\
      \textbf{Stato}        & \stato                 \\
      \textbf{Uso}          & \uso                   \\
      \textbf{Destinato a}  & Imola Informatica      \\
                            & GruppOne               \\
                            & Prof. Tullio Vardanega \\
                            & Prof. Riccardo Cardin  \\
    \end{tabular}
  \end{table}

  \vspace{3cm}
  \textbf{Descrizione}\\
  \descrizione\\
  \vfill
  \verb|gruppone.swe@gmail.com|
\end{center}

\newpage
\thispagestyle{nopage}

\section*{Registro delle modifiche}
\label{sec:registro_delle_modifiche}

\begin{table}[H]
  \label{tab:registro_delle_modifiche}

  \centering
  \rowcolors{2}{lightgray}{white!80!lightgray!100}

  \begin{longtable}[c]{c c c c l}
    \rowcolor{darkgray!90!}\color{white}{\textbf{Versione}} & \color{white}{\textbf{Data}} & \color{white}{\textbf{Nominativo}} & \color{white}{\textbf{Ruolo}} & \color{white}{\textbf{Descrizione}} \\\endhead
    \modifiche
  \end{longtable}
\end{table}

% section registro_delle_modifiche (end)
\newpage

\thispagestyle{nopage}
\pagenumbering{roman}
\tableofcontents

\newpage

\pagenumbering{arabic}


\section{Informazioni logistiche}%
\label{sec:informazioni_logistiche}

\begin{description}
  \item [Luogo] videochiamata su Google Hangouts
  \item [Data] 20/04/2020
  \item [Ora] 18:30 \symbol{8594} 20:00
\end{description}

\subsection{Membri del gruppo presenti}%
\label{sub:membri_del_gruppo_presenti}

\begin{enumerate}
  \item Riccardo Agatea
  \item Tobia Apolloni
  \item Riccardo Cestaro
  \item Alberto Cocco
  \item Luca Ercole
  \item Alberto Gobbo
  \item Alessandro Rizzo
  \item Fabio Scettro
\end{enumerate}
% sub:membri_del_gruppo_presenti (end)

\section{Introduzione}%
\label{sec:introduzione}

L'incontro, che si è tenuto su Google Hangouts, è servito per ripianificare il lavoro.

\section{Ordine del giorno}%
\label{sec:ordine_del_giorno}
\begin{itemize}
  \item Aggiornamento foglio dei requisiti.
  \item Manuale utente e manutentore.
  \item Diagrammi UML\@.
\end{itemize}

\section{Aggiornamento foglio dei requisiti}%
\label{sec:aggiornamento_foglio_dei_requisiti}

Il gruppo ha controllato tutti i requisiti presenti nel foglio Google creato in data 13/04 al fine di verificarne lo stato.
La web-app risulta essere in linea con la pianificazione prevista, mentre app mobile e server sono leggermente indietro.
Si è inoltre provveduto ad assegnare ogni requisito in progress ad un componente del gruppo.

\section{Manuale utente e manutentore}%
\label{sec:manuale_utente_e_manutentore}

Il gruppo ha discusso i due nuovi documenti che è necessario realizzare: manuale utente e manutentore.
Si è deciso di assegnare a Riccardo Cestaro e Alberto Gobbo il compito di definire una struttura per i due documenti.

\section{Diagrammi UML}%
\label{sec:diagrammi_uml}

Il gruppo ha parlato della necessità di avere diagrammi UML corretti per la Product Baseline.
Si è deciso che nell'allegato tecnico saranno consegnati diagrammi delle classi e di sequenza per le tre componenti.

\newpage
\section{Registro delle decisioni}%
\label{sec:registro_delle_decisioni}

\begin{description}
  \item [20200420-int-001] Riccardo Cestaro e Alberto Gobbo si occuperanno di abbozzare la struttura di manuale utente e manutentore.
  \item [2020420-int-002] Nell'allegato tecnico verranno riportati diagrammi delle classi e di sequenza.
\end{description}
% sec:registro_delle_decisioni (end)
\end{document}
